\section{System 0: Architectural requirements for zero-flux substrate}
\label{sec:System0_vacuum_requirements}

System 0 asks: \emph{What geometric properties must a substrate possess to satisfy the three existential requirements of Persistence Mechanics (Finiteness, Conservation, and Causality) when no external memory, boundary conditions, or synchronization signals are available?}

\subsection{Persistence Requires Finiteness for \texorpdfstring{$\mathrm{CAP}, \mathrm{MAR}$}{CAP, MAR}}
\label{sec:finiteness}
The Closure Constraint requires self-specification: it cannot require an external decompressor. Information-theoretically, infinite descriptive complexity implies infinite evaluation time, any finite processing rate yields divergent latency, paralyzing the system before the next fluctuation. Capacity bounds the finite state-space; Margin bounds the reserves against variance. Geometrically, both require the substrate to be discrete, homogeneous, and a lattice.

\subsubsection{Finiteness Mandates a Discrete, Homogeneous Lattice}
\label{subsubsec:discrete_homogeneous_lattice}

A continuous volume admits infinite descriptive complexity. Since there cannot exist an external algorithm (by the Closure Constraint) to compress this to finite Kolmogorov complexity, the substrate must be \textbf{Discrete}: a fundamental, indivisible unit of information setting a maximum density.

Furthermore, to minimize algorithmic processing overhead, the substrate must be \textbf{Homogeneous}. Any spatial variation in local topology would require external memory to decode each neighborhood's structure, violating the Closure Constraint. The local topology must therefore remain constant regardless of location or direction, yielding a constant Kolmogorov complexity independent of system size.

A structure that is both discrete (finite information density) and homogeneous (constant structure in every direction) is by definition a \textbf{Lattice}. This conclusion specifically rules out continuous manifolds (infinite description) and arbitrary graphs (non-constant complexity).

\subsection{Persistence Requires Conservation for \texorpdfstring{$\mathrm{MAP}, \mathrm{PRO}$}{MAP, PRO}}
\label{sec:conservation}

For the vacuum to persist, it must maintain a stable Map over time. This requires that information is perfectly conserved. The vacuum must be a perfectly closed and lossless information network, making Conservation a structural necessity, corresponding to the Map and Protocol pillars. In the quantum vacuum limit, the requirement of Conservation (preserved informational distinguishability) corresponds exactly to quantum unitarity. To maintain substrate independence, we restrict the term Conservation to the architectural requirement, and unitarity exclusively to its macroscopic quantum dual. Geometrically, it imposes four structural properties on the substrate:

\subsubsection{Lattice Must Be Positive Definite}
\label{subsec:positive_definite}
\textit{Information-Theoretic Justification}: In a metric space, the norm $|\mathbf{v}|^2 = g_{\mu\nu}v^\mu v^\nu$ represents the information distance from the origin (reference state). For the system to have a well-defined minimum information state (entropic ground state), this distance must satisfy:
\begin{equation}
    |\mathbf{v}|^2 \geq 0 \quad \forall \mathbf{v} \neq 0
\end{equation}

A indefinite metric with negative eigenvalues (Lorentzian) allows $|\mathbf{v}|^2 < 0$, implying an imaginary ``distance'' from the reference state and preventing the definition of a stable minimum configuration. Furthermore, an indefinite metric (Lorentzian metric) admits non-trivial null vectors ($|\mathbf{v}|^2 = 0$ where $\mathbf{v} \neq 0$), allowing for the creation of infinite information density without exceeding the capacity budget ($|k\mathbf{v}|^2 = 0$ for any scalar amplitude $k$), which violates the \textbf{Finiteness} component. Therefore, the \textbf{substrate} must be Euclidean.

\textbf{The Substrate-Projection Dual:}
The Euclidean positive-definiteness applies strictly to the substrate lattice, which functions as a static memory repository. Causal progression necessitates an active processor operating via an indefinite metric. This fundamental distinction between static storage and dynamic processing provides the geometric degree of freedom from which the directional arrow of time can be derived (\Cref{sec:metric_signature}), satisfying a core causal requirement of the persistence architecture.

\subsubsection{Lattice Must Be Mappable to a Torus}
\label{sec:torus}

To satisfy \textbf{Finiteness}, the system must be spatially bounded. To satisfy \textbf{Protocol} (Conservation), it must be closed (no edges). By the combinatorial Gauss-Bonnet theorem, a regular lattice of fixed vertex coordination tiled on a closed surface of Euler characteristic $\chi$ requires topological defects if $\chi \neq 0$. The sphere ($\chi = 2$) therefore requires defects that break translational invariance. The Euclidean \textbf{Torus} ($T^n$, $\chi=0$) is the unique flat, closed topology with periodic boundary conditions that admits a homogeneous, defect-free lattice structure, simultaneously satisfying Finiteness (bounded), Conservation (closed), and Homogeneity (flat).

Consequently, the statistical evolution of the vacuum is defined by the Partition Function on a torus\footnote{Here, $\beta$ is the formal periodicity parameter of the imaginary time cycle. We use the partition function formalism for state enumeration on a torus, not to claim the vacuum has a literal temperature.}:
\begin{equation}
    Z(\beta) = \sum_{\text{states}} e^{-S_{\text{config}}} = \text{Tr}(e^{-\beta H})
\end{equation}

\subsubsection{Lattice Must Be Even}
\label{sec:even}
A stable Map requires path independence, which requires evenness. If the state count $Z$ depended on the path through moduli space (parameterization) rather than the configuration itself, the vacuum would lack a stable \textbf{Map}.

This requires $Z(\beta)$ to be single-valued under the modular transformation $T: z \to z+1$. The Jacobi Theta Function, $\Theta_\Lambda(z) = \sum e^{i \pi z |\mathbf{v}|^2}$, counts these states. For \textbf{Even} lattices ($|\mathbf{v}|^2 = 2n$), the exponent $2\pi i n z$ is invariant under the shift. However, any odd-norm vector ($|\mathbf{v}|^2 = 2n+1$) introduces a sign inversion ($\Theta \to -\Theta$), rendering the Map multi-valued. Thus, the lattice must be \textbf{Even}.

\subsubsection{Lattice Must Be Self-Dual and Unimodular (Read/Write Symmetry)}
To prevent information loss via destructive interference or Landauer erasure, the encoding operation (write) and the decoding operation (read) must be informationally equivalent. This is enforced by the modular transformation $\mathcal{S}: z \to -1/z$, which maps the lattice to its reciprocal (Fourier dual).

By the Poisson Summation formula, the partition function transforms as:
\begin{equation}
    \Theta_\Lambda(-1/z) \propto \frac{1}{\text{vol}(\Lambda)} \Theta_{\Lambda^*}(z)
\end{equation}
For the Map to remain invariant ($Z_{\Lambda} = Z_{\Lambda^*}$), the lattice must be \textbf{Self-Dual} ($\Lambda = \Lambda^*$). This strictly enforces \textbf{Unimodularity} ($\text{vol}(\Lambda) = 1$), as the only scalar satisfying $V = 1/V$ is $1$. Any other volume introduces an irreversible scaling factor, violating Conservation.

\textbf{Requirements:} The lattice of the vacuum must be \textbf{Positive Definite, Unimodular, Even, and Self-Dual}.

\subsection{Persistence Requires Causality for \texorpdfstring{$\mathrm{GOV}, \mathrm{TOL}$}{GOV, TOL}}
\label{sec:causality}
Persistence requires stable evolution. Information-theoretically, this is the Serialization Problem: projecting a finite state-space onto a discrete timeline. Geometrically, this requires a well-defined temporal axis without ambiguity.

\textbf{The Serialization Problem:} From the \textbf{Capacity} component, any persistent system possesses a finite state-space dimension $N$ (the number of linearly independent, orthogonal basis vectors spanning the node's total informational payload, i.e., the Hilbert-space dimension). To evolve, the system must serialize these $N$ states onto a discrete timeline of $\Delta$ slots per fundamental cycle, the \textbf{Temporal Modulus}.

\textbf{Causal Aliasing:} By the Pigeonhole Principle, if $N > \Delta$, at least one temporal slot must host multiple distinct states. This \textbf{Causal Aliasing} destroys the system's ability to maintain a well-defined serial history.
    
\textbf{The Persistence Inequality:} To enforce a strict arrow of time and satisfy the Causality requirement, the projection must satisfy $\Delta \ge N$. The temporal container ($\Delta$, the Toll) must meet or exceed the state content it constrains ($N$, via the Governor).

\subsection{Summary: The Architectural Requirement of the Vacuum}
The architectural requirements for the substrate are:

\begin{enumerate}
    \item ($\mathrm{CAP}, \mathrm{MAR}$) \textbf{Finite Lattice}, required by Finiteness.
    \item ($\mathrm{MAP}, \mathrm{PRO}$) \textbf{Positive Definite, Unimodular, Even, and Self-Dual}, required by Conservation. 
    \item ($\mathrm{GOV}, \mathrm{TOL}$) Any causal system must satisfy a \textbf{Causal Projection}: a serialization of its $N$-dimensional state-space onto a discrete timeline of length $\Delta \ge N$.
\end{enumerate}

With this specification complete, identifying the physical substrate becomes a strict eliminative test: to determine whether a unique, self-specifying geometry exists that satisfies all of these requirements simultaneously.